\chapter{Triangles isocèles}
\begin{itemize}
\it Soit un triangle $ABC$ isocèle en $A$. On prolonge $[BA)$ d'une longueur $AD=BA$. \begin{enumerate}
\item Montrer que le triangle $ACD$ est isocèle. 
\item Montrer que l'un des angles du triangle $BCD$ est égal à la somme des deux autres. 
\item Si dans un triangle $ABC$, l'angle en $A$ est égal à la somme des angles en $B$ et $C$, montrer que la médiane $AM$ est égale à la moitié du côté $[BC]$.
\end{enumerate}
\it On construit un cercle de centre $O$ et deux diamètres $[AC]$ et $[BD]$.\begin{enumerate}
\item Démontrer que le quadrilatère $ABCD$ a ses quatres angles égaux. 
\item Quels sont les axes de symétrie de ce quadrilatère ? En déduire l'égalité de $AB$ et $CD$, ainsi que de $AD$ et $BC$.
\end{enumerate}
\it Dans un triangle quelconque $ABC$, la bissectrice extérieure de l'angle en $A$ coupe en $D$ le prolongement $[Cx)$ du côté $[BC]$. On prolonge $[BA]$ d'une longueur $AE=AC$. \begin{enumerate}
\item Que représente la droite $(DA)$ pour le segment $[CE]$, pour le triangle $BDE$ ? 
\item Calculer la mesure de l'angle $BED$ sachant que $\widehat{ACB}=84^o$.
\end{enumerate}
\it On prolonge d'une même longueur et dans le même sens de parcours, les trois côtés d'un triangle équilatéral $ABC$. On obtient les points $A'$ sur le prolongement de $[BC]$, $B'$ sur celui de $[CA]$, et $C'$ sur celui de $[AB]$.\begin{enumerate}
\item Montrer que l'on peut par glissement amener le 
calque de la figure obtenue $ABCA'B'C'$ sur $BCAB'C'A'$ ou $CABC'A'B'$.
\item En déduire que sur les trois triangles $A'B'C$, $B'C'A$ et $C'A'B$ sont directement égaux. Quelle est la nature du triangle $A'B'C'$ ?
\end{enumerate}
\it Soit un quadrilatère convexe $ABCD$ dans lequel $AB=BC$ et $\widehat{BAD}=\widehat{BCD}$.
\begin{enumerate}
\item Montrer que les angles $\widehat{DAC}$ et $\widehat{DCA}$ sont égaux. Comparer $DA$ et $DC$.
\item Que représente $(BD)$ pour le segment $[AC]$ et pour les angles $\widehat{ABC}$ et $\widehat{ADC}$ ?
\end{enumerate}
\it Démontrer que si dans un triangle $ABC$ l'une des conditions suivantes est remplie, le triangle est isocèle. \begin{enumerate}
\item La hauteur $(AH)$ est en même temps bissectrice de l'angle $\widehat{A}$.
\item La hauteur $(AH)$ est en même temps médiane relative à $[BC]$. 
\item La médiatrice de $[BC]$ passe par le sommet $A$.
\end{enumerate}
\it Soit un triangle $ABC$ dans lequel la médiane $(AM)$ est bissectrice de l'angle intérieur $\widehat{BAC}$. On prolonge $[AM]$ d'une longueur $MD=AM$.\begin{enumerate}
\item Montrer que les triangles $MAB$ et $MDC$ sont symétriques par rapport à $M$. En déduire que les angles $\widehat{MAB}$ et $\widehat{MDC}$ sont égaux ainsi que les segments $[AB]$ et $[DC]$. 
\item Nature du triangle $CAD$ ? Comparer $AB$ et $AC$.
\item Énoncer le théorème qui en résulte.
\end{enumerate}
\it Dans un triangle isocèle $ABC$ de base $[BC]$, on mène la médiane $(BM)$. Soit $N$ le milieu de $[AB]$ et $G$ le point de rencontre de $[BM]$ avec la hauteur $(AH)$.\begin{enumerate}
\item On retourne le triangle $ABC$ sur lui-même en amenant $B'$ calque de $B$ en $C$ et $C'$ calque de $C$ en $B$. Où viennent les calques $M'$ et $N'$ de $M$ et de $N$ ? 
\item Montrer que les points $C$, $G$ et $N$ sont alignés et que les triangles $GBC$ et $GMN$ sont isocèles. Comparer $BM$ et $CN$. 
\item Que représente $(AH)$ pour le segment $[MN]$ ?
\end{enumerate}
\it On considère un triangle $ABC$ isocèle en $A$. La bissectrice intérieure de l'angle en $B$ coupe le côté $[AC]$ en $D$ et la hauteur $(AH)$ en $I$. \begin{enumerate}
\item Nature du triangle $IBC$ ? Montrer que $[CI)$ est la bissectrice intérieure de l'angle en $C$. 
\item La bissectrice $[CI)$ coupe le côté $[AB]$ en $E$. Montrer que le symétrique de $D$ par rapport à la droite $(AH)$ est le point $E$. Comparer $AD$ et $AE$, puis $ID$ et $IE$. 
\item Démontrer que $BD=CE$.
\end{enumerate}
\it on considère un triangle $AHB$ rectangle en $H$. 
\begin{enumerate}
\item On suppose que $AB=2HB$. Démontrer que l'angle aigu en $B$ est le double de l'angle aigu en $A$ (prolonger $[BH]$ de $HC=BH$).
\item On suppose que par hypothèse l'angle en $B$ soit le double de l'angle en $A$. Démontrer que l'on a $AB=2HB$.
\end{enumerate}
\it \begin{enumerate}
\item Construire un quadrilatère convexe $ABCD$ dans lequel les trois angles en $D$, $A$, et $B$ sont égaux à $100^o$ et les côtés $[AB]$ et $[AD]$ égaux à $20$ mm.
\item Démontrer que le triangle $BCD$ est isocèle et comparer $CB$ et $CD$.
\item Que représente la droite $(AC)$ pour le segment $[BD]$ et pour les angles $\widehat{DAB}$ et $\widehat{BCD}$ ?
\end{enumerate}
\it On considère un angle aigu $\widehat{xOy}$ et un point intérieur $A$. \begin{enumerate}
\item Construire les symétriques $B$ et $C$ du point $A$ par rapport à $(Ox)$ et $(Oy)$. Montrer que les points $A$, $B$ et $C$ sont sur un cercle de centre $O$.
\item On mène la droite $(BC)$ qui coupe $(Ox)$ en $M$
et $(Oy)$ en $N$. Comparer les angles $\widehat{OBC}$
et $\widehat{OCB}$, puis montrer que $(AO)$ est la bissectrice de l'angle $\widehat{MAN}$.
\item Soit $D$ le symétrique de $C$ par rapport à $(Ox)$. Démontrer que les trois points $A$, $M$, $D$ sont alignés.
\end{enumerate}



\end{itemize}